\subsection{本原元素定理}

设 E 是 K 的一个扩张；E 中的元素 x 称为本原的，如果 \(E = K[x]\) 。要使扩张 E 具有本原元素，必须 \([E : K]\) 是有限的。

定理1。—— 设 E 是 K 的一个扩张。则以下条件等价：

\begin{enumerate}
\def\labelenumi{\alph{enumi})}
\item
  E 具有本原元素；
\item
  E 的子扩张只有有限多个。
\end{enumerate}

当 E 是有限次可分扩张时，这些条件成立。

首先假设 E 具有本原元素 x，并设 f 是 x 在 K 上的极小多项式。对每个首一多项式 \(g \in E[X]\) （在 E{[}X{]} 中整除 f 的），用 \(E_g\) 表示由 g 的系数生成的 E 的子扩张。由于可能的多项式 g 数量有限（若 f 在 E{[}X{]} 中分解为 r 个首一不可约多项式的乘积，则该数量以 \(2^{r}\) 为界），子扩张 \(E_g\) 只有有限多个。因此为证明 b)，只需证明 E 的每个子扩张 L 都是这些 \(E_g\) 之一。现在若 L 是 E 的一个子扩张，则我们有 L{[}x{]} = E；若 g 是 x 在 L 上的极小多项式，我们有 {[}E:L{]} = deg(g)。此外，g 是 f 在 L{[}X{]} 中的因子，从而也在 E{[}X{]} 中；所以我们有 \(E_g \subset L\) 且 \(E = E_g[x]\) 。由于 \(g(x) = 0\) ，我们有 \([E : E_g] \leq \deg(g)\) ，因此 \([E : E_g] \leq [E : L]\) ，从而 \(L = E_g\) ，此即所要证明的。

接下来我们观察到，条件 b) 蕴含扩张 E 是有限次的：根据 V 第 18 页的注记 2，只需证明它是代数的；现在若 z 是 E 中在 K 上超越的元素，则子扩张 \(\mathbf{K}(z^{n})\) , \(n \in N\) 两两不同。

为了证明 \(b) \Rightarrow a)\) 我们现在区分两种情况：

\begin{enumerate}
\def\labelenumi{\Alph{enumi})}
\item
  若域 K 是有限的，则域 E 是 K 上的有限维向量空间，从而是有限集。所以 \({}^{1}\) （V，第78页，引理1）存在 E 的一个元素 x 生成 E 的乘法群，并且我们有 \(E = K[x]\) 。
\item
  现在假设域K是无限的。若b)成立，则扩张E是有限次的，因此b)也可以表述为：E只含有有限多个子代数。既然如此，蕴含关系 \(b) \Rightarrow a)\) 是以下更一般命题的推论（其中域K为无限的假设是不可或缺的，参见V，第153页，§7的习题5）：
\end{enumerate}

命题7. --- 假设 K 是无限的；设 A 是一个交换 K-代数，它只拥有有限多个子代数（例如一个平展 K-代数，V，第30页，命题3），并设 V 是 A 的一个生成 A 的向量子空间。则存在 \(x \in V\) 使得 \(A = K[x]\) 。

设 \(A_{1}, \ldots, A_{n}\) 为 A 的异于 A 的子代数。若 \(x \notin A_{1} \cup \ldots \cup A_{n}\) ，则子代数 \(K[x]\) 不能等于任何一个 \(A_{i}\) ，从而必与 A 重合。进而，由于 V 生成 A，它不包含于子空间 \(A_{i}\) 中的任何一个。因此命题7是以下引理的推论：

引理1. --- 设 A 为向量 K-空间，V 与 \(\mathbf{A}_1, \ldots, \mathbf{A}_n\) 为 A 的子空间。如果 \(\operatorname{Card}(\mathbf{K}) \geqslant n\) 且 V 不包含于任何 \(\mathbf{A}_i\) ，则 V 不包含于 \(\mathbf{A}_1 \cup \ldots \cup \mathbf{A}_n\) 。

对 \(n\) 进行归纳论证，我们只需证明：若 \(\mathbf{V} \subset \mathbf{A}_n\) 且 \(\mathbf{V} \subset \mathbf{A}_1 \cup \ldots \cup \mathbf{A}_n\) ，则 \(\mathbf{V} \subset \mathbf{A}_1 \cup \ldots \cup \mathbf{A}_{n-1}\) 。设 \(x \in \mathbf{V}\) , \(x \notin \mathbf{A}_n\) ，并设 \(y\) 为 \(\mathbf{V}\) 中的任意元素。若 \(y \in \mathbf{K}x\) ，我们有 \(y \in A_1 \cup \ldots \cup A_{n-1}\) ；若不然，则元素 \(x\) 与 \(y + \lambda x\) （\(\lambda \in \mathbf{K}\)）的个数严格大于 \(n\) ，且都属于 \(\mathbf{A}_1 \cup \ldots \cup \mathbf{A}_n\) ，因此其中有两个属于同一个 \(\mathbf{A}_i\) 。于是存在 \(i\) , \(1 \leqslant i \leqslant n\) ，使得要么 \(x \in \mathbf{A}_i\) 且 \(y + \lambda x \in \mathbf{A}_i\) （对某个 \(\lambda \in \mathbf{K}\)），要么 \(y + \lambda x \in \mathbf{A}_i\) 且 \(y + \mu x \in \mathbf{A}_i\) （对两个不同的标量 \(\lambda, \mu \in \mathbf{K}\)）。在两种情形中我们都得出 \(x \in \mathbf{A}_i\) 且 \(y \in \mathbf{A}_i\) ；但这意味着 \(i \neq n\) ，因此 \(y \in \mathbf{A}_1 \cup \ldots \cup A_{n-1}\) ，此即所要证明的。

这就完成了定理1中 \(a)\) 与 \(b)\) 等价性的证明。最后，若扩张 \(E\) 是可分的且有限次的，则条件 \(b)\) 成立，由 \(V\) 第30页命题3。

